\begin{document}
\chapter{Introduction}
The Problem: Malaria diagnosis usually relies on manual microscopy (a human looking through a microscope). This is slow, labor-intensive, and prone to human error (fatigue, lack of expertise).

The Goal: Automate the detection of parasitized cells to assist doctors, speed up diagnosis, and improve accuracy in resource-poor regions.

Quantity (#): Total of 27,558 cell images.

Quality/Balance: The dataset is perfectly balanced (50% Parasitized, 50% Uninfected), which is excellent for training.

Supervision: This is a Supervised Learning problem (the data is labeled).

Input Data: The images are RGB (color) but vary in resolution/size, meaning preprocessing is required.

\chapter{State of the Art}

Traditional Methods: Historically, people used manual feature extraction (HOG, SIFT) combined with SVM classifiers.

Modern Methods (Deep Learning): The standard solution now is Convolutional Neural Networks (CNNs).

Specific Architectures: Literature shows high success rates using Transfer Learning with models like VGG-16, ResNet50, or DenseNet.

Comparative Analysis: "I want to compare a custom-built lightweight CNN against a heavy Transfer Learning model to see which is more efficient."

Metric Focus: "Most models focus on Accuracy; my goal is to maximize Recall (minimizing False Negatives) because missing a sick patient is dangerous."

Explainability: "I want to not just classify, but use Grad-CAM to show where the parasite is located in the image."

\end{document}